\documentclass[11pt,a4paper]{article}
\usepackage[margin=25mm]{geometry}
\usepackage{amsmath,amssymb,bm,mathtools}
\usepackage[T1]{fontenc}
\usepackage{lmodern,microtype}
\usepackage[colorlinks=true,linkcolor=blue,urlcolor=blue]{hyperref}
\newcommand{\R}{\mathbb R}
\title{The scalar reduction behind \texttt{eq:linear-soln}:\\the power of $c$, the definition of $h_q$, and the corrected equation}
\author{}
\date{}
\begin{document}
\maketitle

\section{Setting}
Section~2 studies a linear network with one hidden layer of $N$ spatially embedded units, spatial kernel $K = F\Lambda F^\dagger$ (real, symmetric, translation invariant, eigenvalues $\lambda_q = \lambda_{-q}$), input and readout priors $k_0 I_m$ and $k_1 I_n$, rate penalty $g^2 I$, input noise $\sigma^2$ and $\kappa = \beta/(Pn)$.
Under the shared-basis ansatz ($XX^\top = U S_X^2 U^\top$, $YY^\top = U S_Y^2 U^\top$, $\Sigma_q = U D_q U^\top$, which requires $[XX^\top, YY^\top] = 0$) every matrix is diagonal in $U$ and the problem decouples into one scalar problem per data mode $p = 1,\dots,P$.
For one such mode write $s_X^2, s_Y^2$ for the diagonal elements of $S_X^2, S_Y^2$, and $a_q, b_q, d_q, c$ for those of $D_{A,q}, D_{B,q}, D_q, D_C$. With
\begin{equation}
  \hat s_X^2 = \frac{k_0 s_X^2}{m}, \qquad \hat s_Y^2 = \frac{s_Y^2}{k_1}, \qquad \hat\sigma_q^2 = \frac{\sigma^2}{\lambda_q}, \qquad \delta = \frac{1}{k_1\kappa},
\end{equation}
the matrix equations of Section~2 are
\begin{equation}
  D_{A,q} = \lambda_q\tfrac{k_0}{m}S_X^2 + \sigma^2 I,\qquad
  D_{B,q} = \frac{\lambda_q}{k_1}D_C^{-2}S_Y^2 - g^2 I,\qquad
  D_q = \big(D_{A,q}^{-1} - D_{B,q}\big)^{-1},\qquad
  D_C = \frac1N\sum_q\lambda_q D_q + \delta I .
  \label{eq:matrix}
\end{equation}
(The $1/N$ in $D_C$ is the subject of the companion note \texttt{rg\_kernel\_normalisation.tex}.)

\section{\texorpdfstring{Problem 1: the power of $c$ in $b_q$}{Problem 1: the power of c in b\_q}}
$\tilde B_q = \frac{\lambda_q}{k_1}\tilde C^{-1}YY^\top\tilde C^{-1} - g^2 I$ is \emph{quadratic} in $\tilde C^{-1}$, so in the shared basis $D_{B,q}$ contains $D_C^{-2}$, as Section~2 itself writes in \eqref{eq:matrix}. Its diagonal element is therefore
\begin{equation}
  b_q = \frac{\lambda_q s_Y^2}{k_1 c^2} - g^2 = \frac{\lambda_q\hat s_Y^2}{c^2} - g^2 ,
\end{equation}
whereas the scalar reduction in Section~2 used $\lambda_q\hat s_Y^2/c$.
A units check confirms the square: $\tilde C$ has the units of a kernel, $[\tilde C] = [\phi^2]$, and the tilt $b_q$ must have units $[\phi^{-2}]$ to sit next to $a_q^{-1}$ in $d_q^{-1} = a_q^{-1} - b_q$; with $[\hat s_Y^2] = [\phi^2]$ (since $k_1\tilde C$ and $YY^\top$ are both output kernels), only $\hat s_Y^2/c^2$ has the right units.

\section{\texorpdfstring{Problem 2: the definition of $h_q$}{Problem 2: the definition of h\_q}}
Section~2 defines $h_q$ through
\begin{equation}
  a_q^{-1} + g^2 = \frac{1}{\lambda_q\hat s_X^2 + \sigma^2} + g^2 \overset{!}{=} \frac{\lambda_q}{h_q}.
\end{equation}
Solving,
\begin{equation}
  h_q = \frac{\lambda_q(\lambda_q\hat s_X^2 + \sigma^2)}{1 + g^2(\lambda_q\hat s_X^2+\sigma^2)}
      = \frac{\lambda_q\hat s_X^2 + \sigma^2}{\lambda_q^{-1} + g^2\big(\hat s_X^2 + \hat\sigma_q^2\big)} .
  \label{eq:hq}
\end{equation}
Section~2 had $g^2(\lambda_q\hat s_X^2 + \hat\sigma_q^2)$ in the denominator: one factor of $\lambda_q$ too many on $\hat s_X^2$ (and its $\hat\sigma_q^2$ was defined circularly as $\hat\sigma^2/\lambda_q$; it must be $\sigma^2/\lambda_q$).
Multiplying out shows the inconsistency: the version in Section~2 gives
$\lambda_q/h_q = \big(1 + g^2\lambda_q^2\hat s_X^2 + g^2\sigma^2\big)/(\lambda_q\hat s_X^2+\sigma^2)$, which differs from $a_q^{-1}+g^2$ by $g^2(\lambda_q^2-\lambda_q)\hat s_X^2/(\lambda_q\hat s_X^2+\sigma^2)$, nonzero whenever $g\neq0$ and $\lambda_q\neq1$.
At $g = 0$ both reduce to $h_q = \lambda_q(\lambda_q\hat s_X^2+\sigma^2) = \lambda_q a_q$, so the error is invisible in the $g\to0$ limit but changes every result at finite rate penalty. For example, at large $g$ the correct $h_q \to \lambda_q/g^2$, whereas the incorrect one tends to $\lambda_q(\lambda_q\hat s_X^2+\sigma^2)/\big(g^2(\lambda_q^2\hat s_X^2+\sigma^2)\big)$, which retains a spurious dependence on the input spectrum.

\section{Derivation of the corrected implicit solution}
With both corrections,
\begin{equation}
  d_q^{-1} = a_q^{-1} - b_q = \frac{\lambda_q}{h_q} - \frac{\lambda_q\hat s_Y^2}{c^2}
  \quad\Longrightarrow\quad
  d_q = \frac{c^2\,h_q/\lambda_q}{c^2 - \hat s_Y^2 h_q}.
\end{equation}
Positivity of the precision, $d_q > 0$ for every $q$, requires $c^2 > \hat s_Y^2\max_q h_q$.
The self-consistency for $c$ (with the $1/N$ of the companion note) is
\begin{equation}
  c = \frac1N\sum_q\lambda_q d_q + \delta = \frac{c^2}{N}\sum_q\frac{h_q}{c^2-\hat s_Y^2h_q} + \delta ,
\end{equation}
and dividing by $c$ gives the corrected \texttt{eq:linear-soln}:
\begin{equation}
  \boxed{\;1 - \frac{\delta}{c} = \frac{c}{N}\sum_q\frac{h_q}{c^2 - \hat s_Y^2h_q},
  \qquad
  h_q = \frac{\lambda_q(\lambda_q\hat s_X^2 + \sigma^2)}{1 + g^2(\lambda_q\hat s_X^2+\sigma^2)}\;}
  \label{eq:linear-soln-corrected}
\end{equation}
to be solved for $c > \hat s_Y\sqrt{\max_q h_q}$, separately for each data mode. Equivalent forms that are sometimes more convenient are
\begin{equation}
  c - \delta = \frac1N\sum_q \frac{c^2 h_q}{c^2 - \hat s_Y^2 h_q}
  \qquad\text{and}\qquad
  c - \delta - \frac1N\sum_q h_q = \frac{\hat s_Y^2}{N}\sum_q\frac{h_q^2}{c^2 - \hat s_Y^2h_q},
\end{equation}
the second of which isolates the learning term: at $\hat s_Y = 0$ it gives $c = \delta + \frac1N\sum_q h_q$, i.e.\ $d_q = h_q/\lambda_q$, the prior covariance (shrunk by the rate penalty).
As $N\to\infty$ at fixed lattice spacing, $\frac1N\sum_q\to\int_{\mathrm{BZ}}\frac{\mathrm d^dq}{(2\pi)^d}$.

\section{Consequences}
\begin{itemize}
  \item The structure of the solution changes. In the original equation $1-\delta/c = \sum_q h_q/(c-\hat s_Y^2h_q)$ each summand is invariant under a joint rescaling of $c$ and $h_q$; in the corrected one it is not, so the balance between the target strength $\hat s_Y^2$ and the spatial spectrum $h_q$ is different. The singular threshold moves from $c = \hat s_Y^2\max_q h_q$ to $c = \hat s_Y\sqrt{\max_q h_q}$. Any threshold, onset of mode-selective learning or scaling law read off from the original equation has to be re-derived.
  \item The finite-$g$ behaviour changes through $h_q$, in particular the large-$g$ limit and how the rate penalty shapes the spatial spectrum $h_q$.
  \item Any later analysis (limiting cases, continuum limit, numerical predictions) must use \eqref{eq:linear-soln-corrected}.
\end{itemize}

\paragraph{Numerical check.}
On a ring of $N=12$ sites with a Gaussian circulant $K$, $P=4$ patterns, $X = ZW_X^\top$, $Y = ZW_Y^\top$ (so the data kernels commute), and $k_0 = 0.7$, $k_1 = 1.3$, $\sigma^2 = 0.2$, $g^2=0.3$, $\delta = 0.5$, iterating the full $NP\times NP$ Gaussian fixed point (no ansatz, no scalar reduction) and diagonalising $C + \delta I$ gives
$c \in \{0.6887, 0.7020, 1.1059, 2.4602\}$; the roots of \eqref{eq:linear-soln-corrected} agree to all printed digits.
The equation as originally written gives $\{2.764, 2.884, 5.168, 9.700\}$.

\section{Recommended fix}
Use $b_q = \lambda_q\hat s_Y^2/c^2 - g^2$, define $h_q$ by \eqref{eq:hq} with $\hat\sigma_q^2 = \sigma^2/\lambda_q$, carry $c^2$ through $d_q$ and the self-consistency, and replace \texttt{eq:linear-soln} by \eqref{eq:linear-soln-corrected}. These edits are marked in red in Section~2, with the label \texttt{eq:linear-soln} kept on the corrected equation.
\end{document}
